% Compile twice: lualatex 20001182.tex
\documentclass[11pt]{article}
\usepackage[margin=25mm]{geometry}
\usepackage{fontspec,unicode-math}
\setmainfont{Latin Modern Roman}
\setmathfont{Latin Modern Math}
\usepackage{amsmath,amssymb,mathtools}
\usepackage{xurl,hyperref}
\hypersetup{colorlinks=true,urlcolor=blue,linkcolor=blue}
\setcounter{secnumdepth}{0}
\setlength{\parindent}{0pt}
\setlength{\parskip}{5pt}
\setlength{\emergencystretch}{3em}
\newcommand{\sref}[2]{\hyperlink{source:#1:#2}{[S#2]}}
\newcommand{\cataloguescope}[1]{\paragraph{Source question.}#1}
\newcommand{\floor}[1]{\left\lfloor#1\right\rfloor}
\begin{document}
% UnsolvedMath ID 20001182; source catalogue code AIM-COMPUTATION-0020
\setcounter{section}{20001181}
\section{Compatible and complete odds-ratio specifications}\label{20001182}
{\small\sffamily UnsolvedMath ID 20001182 · Computer Science}
\subsection{Definitions and mathematical statement}

\paragraph{Shared notation.}Standard number systems are denoted by $\mathbb N,\mathbb Z,\mathbb Q,\mathbb R,\mathbb C$; index ranges, computational models and specialized notation follow the cited problem source and the definitions below.


A multiway table is a nonnegative array indexed by the levels of several categorical variables. An odds-ratio specification fixes selected multiplicative ratios of its cells. A complete specification identifies a distribution within the stated normalization and margins, while a partial specification leaves a nontrivial family of compatible tables.
\paragraph{Statement recorded in the supplied source.}
Question 6. For multi-way tables, what are the complete and partial specifications that come from specifying sets of odds-ratios? \sref{20001182}{2}
\paragraph{Scope qualification.}
The supplied extraction includes adjacent workshop exposition; the numbered target passage is retained verbatim and that surrounding exposition is omitted. The source asks for characterization across multiway-table models and does not fix a particular dimension or set of ratios.
\paragraph{Source status.}
The supplied catalogue labels this record partially solved. Its source-status metadata was checked against the supplied snapshot on 2026-10-11. This import does not independently establish current open status or mathematical correctness.
\subsection{Short English statement}
Which collections of odds ratios determine a categorical distribution completely, and what do incomplete collections determine?
\subsection{Sources}
\begin{itemize}
\item[\hypertarget{source:20001182:1}{S1}] Ulam AI and the UnsolvedMath Contributors, UnsolvedMath: A Curated Collection of Open Mathematics Problems, record 20001182 (AIM-COMPUTATION-0020). Catalogue curation and metadata: CC BY 4.0. The supplied source question is retained; definitions and the English summary are editorial additions.
\url{https://huggingface.co/datasets/ulamai/UnsolvedMath}.
\item[\hypertarget{source:20001182:2}{S2}] AIM Workshop Problem Lists, original problem source linked by catalogue record AIM-COMPUTATION-0020
\url{https://aimath.org/WWN/compalgstat/compalgstat.pdf}.
\end{itemize}
\paragraph{Preparation and review.}
Original question attributed to the cited source and UnsolvedMath contributors. Codex (GPT-6; exact serving revision not exposed) assisted with definitions, English summaries, source inspection and formatting. No new human mathematical authorship or independent proof verification is asserted. Canonical ID 20001182 and source code AIM-COMPUTATION-0020 checked against the complete supplied registry and existing repository statements; known flagged repetitions and obvious identical conjectures excluded. Full mathematical-equivalence/literature review has not been performed. Statement preparation and final review followed contribution prompts C and D; automated checks verify structure and identity, not open status.
\label{end:20001182}
\end{document}
